\chapter{Durchführung} \label{ch:Durchfuehrung}

\section{Versuchsaufbau/Geräte} \label{sec:Aufbau}

Der Versuchsaufbau besteht aus einem \(\beta\)-Typ Stirlingmotor, der je nach Versuchsteil mit unterschiedlichen Zylinderköpfen betrieben wird. Für den Betrieb als Kältemaschine wird ein Teflonkopf mit integrierter Heizwendel verwendet, während für den Betrieb als Wärmekraftmaschine ein Keramikkopf mit Heizwendel eingesetzt wird. 

Zur Messung der relevanten Größen stehen ein Sensorsystem (Cassy) mit Druck- und Temperatursensor, ein Dreikanalthermometer, ein Multimeter, ein optisches Drehzahlmessgerät sowie ein Durchflussmengenmesser für das Kühlwasser zur Verfügung. Ein regelbares Netzteil versorgt sowohl die Heizwendel als auch den Elektromotor, der den Stirlingmotor im Kältemaschinenbetrieb antreibt. Für die Belastungsmessung wird ein Prony-Zaum mit Federkraftmesser verwendet.

\section{Messverfahren} \label{sec:Messverfahren}

\subsection{Betrieb als Kältemaschine}

Zunächst wird der Stirlingmotor über den Elektromotor angetrieben. Die Temperatur im oberen Zylinderbereich sinkt dabei ab. Durch die Heizwendel wird so lange Wärme zugeführt, bis die Temperatur wieder dem Ausgangswert entspricht. In diesem stationären Zustand entspricht die elektrische Heizleistung der Kälteleistung des Motors. Es werden Heizstrom, Heizspannung, die Temperaturdifferenz des Kühlwassers sowie der Volumenstrom des Kühlwassers aufgenommen.

Im zweiten Schritt wird der Zylinderkopf mit Reagenzglas montiert und etwa \(1\,\mathrm{ml}\) Wasser eingefüllt. Während der Motor als Kältemaschine läuft, wird der Temperaturverlauf des Wassers aufgezeichnet, bis ein stabiles Plateau unterhalb des Gefrierpunkts erreicht ist. 

\subsection{Betrieb als Wärmepumpe}

Ohne Änderung der Motoreinstellungen wird der Motor anschließend im Wärmepumpenmodus betrieben. Die Temperaturverläufe des Kühlwassers und der Wasserprobe werden weiter aufgezeichnet, um den Wärmetransport qualitativ zu analysieren.

\subsection{Betrieb als Wärmekraftmaschine (Leerlauf)}

Für den Betrieb als Wärmekraftmaschine wird der Keramikzylinderkopf montiert und der Motor durch die Heizwendel aufgeheizt. Nach einer kurzen Einschwingzeit läuft der Motor selbstständig. Es werden Heizstrom, Heizspannung, die Drehzahl, die Temperaturdifferenz des Kühlwassers sowie der Volumenstrom aufgenommen. Zusätzlich werden mehrere \(pV\)-Diagramme aufgezeichnet, deren eingeschlossene Fläche die pro Zyklus verrichtete mechanische Arbeit liefert.

\subsection{Belasteter Betrieb mit Prony-Zaum}

Zur Bestimmung der effektiven mechanischen Arbeit wird der Prony-Zaum auf die Motorachse gesetzt und mit einem Federkraftmesser belastet. Für verschiedene Bremskräfte werden jeweils die Drehzahl und die Fläche des \(pV\)-Diagramms gemessen. Aus Kraft und Hebellänge wird das Drehmoment bestimmt, woraus sich die pro Zyklus abgegebene Arbeit ergibt. Diese Messungen ermöglichen die Bestimmung des effektiven Wirkungsgrads des Motors unter Last.
